\documentclass[a4paper,11pt]{article}

\usepackage{revy}
\usepackage{amssymb}
\usepackage{amsmath}
\usepackage{commath}
\usepackage{amsfonts}

\revyname{MatematikRevyen}
\revyyear{2025}
\version{0.1}
\eta{5 minutter}
\status{Færdig}

\title{Bogklubben}
\author{Augusta '20, Frigg '20, Hugo '21}

\begin{document}
\maketitle

\begin{roles}
\role{I}[Sommer] Instruktør
\role{K}[Markus] Kalkulus (Med norsk accent)
\role{D}[Daniel] Dummit and Foote  
\role{L}[Emil] LIM bog 
\role{B}[Lise] Den Lille Blå
\role{A}[Anna] Anna (Ikke på scenen)
\end{roles}

\begin{props}
\prop{1} 4 kæmpe bøger som man kan have på
\end{props}


\begin{sketch}

\scene Lyset går op og Kalkulus, Dummit og LIM står midt på scenen. 

\says{D} Hej allesammen og velkommen tilbage i bogklubben! I denne uge skulle vi alle sammen læses af Anna! Jeg glæder mig meget til at høre jeres betragninger om hende, og håber at vi kan have en spændende diskussion. 

\says{L} Årh Anna! Hun er så sød...

\says{K} Og så smuk...

\scene LIM rækker hånden op.

\says{D} LIM, du har et spørgsmål? 

\says{L} Ja. Hvorfor er Kalkulus med i bogklubben? Han er jo en førsteårs bog. 

\says{D} Det er fordi at Anna dumpede MatIntro sidste år og nu tager det igen.

\says{K} Jah, Julie og jeg har hatt mange lange kvelder sammen for tiden. 

\says{L} \act{Trist} Jeg kan ikke sige det samme. Jeg tror sammenlagt at hun har haft mig åben i en time siden blokken startede. 

\says{D} \act{Også trist} Samme her. Og den time brugte hun på at læse den samme linje igen og igen. 

\says{K} Nåå, ser man det! Nei Anna og jeg hat utviklet et litt spesielt bånd. Sidste søndag tilbrakte Julie og jeg HELE dagen sammen, det var bare SÅ hyggeligt! ...og vil I høre en hemmelighet? 

\scene LIM og Dummit rykker helt ind til Kalkulus. 

\says{K} Det er litt dårligt hihi.. men jeg håper faktisk litt på at hun ender med å skulle ta MatIntro en tredie gang næste år. 

\says{D} Det er i hvert fald helt sikkert at LIM og jeg kommer til at se mere til hende næste år, hvis hun fortsætter på at bruge så lidt tid på os som nu.. 

\scene Den Lille Blå kommer gående ind og virker udmattet. 

\says{B} Hej. Undskyld jeg kommer for sent. 

\says{L} Hvad?! Er han OGSÅ med? 

\says{D} Lille Blå, hvorfor kommer du så meget for sent? 

\says{B} Jeg sov over mig. Anna ville ikke give slip på mig hele natten! Jeg tror hun havde en aflevering for eller noget. 

\says{K} Hvad?! Hun fortalte meg at hun hadde fri i kveld!

\says{B} \act{Griner lidt af K, har næsten ondt af K} Nå, men hun var altså sammen med mig! 

\says{K} Hvad med sidste weekend? Da hun skulle "i teatret med sin mor"? 

\says{B} Troede du virkelig hun var i teateret en hel weekend? Nej, der lavede hun også matematik med mig... 

\says{K} What, hvorfor skulle hun lyve om det? Og hvorfor deg?? Alt hvad Anna kan læse i deg, kan hun da også læse i meg! 

\says{B} Jeg tror måske bare godt at hun kan lide at jeg svarer lidt mere kort og præcist på hendes spørgsmål.

\says{K} Siger du jeg snakker for mye?! 

\says{B} Nej nej... jeg er jo også bare en del lettere at have med i tasken-

\says{K} Kaller du meg fed?!

\says{B} Altså jeg siger bare, af en paperback at være, er du rimelig tung!

\scene Kalkulus og Lille Blå begynder at fronte hinanden meget. 

\says{D} Så så, lad os lige holde den gode tone her! Det hedder bogklubben og ikke fight klubben ahah... \act{kigger lidt håbende ud mod publikum}

\scene Kalkulus og Lille Blå stopper lidt op og kigger på D som den store idiot D er. LIM bogen klemmer sig dramatisk ind imellem Kalkulus og Lille Blå. 

\says{L} Ej venner!! Vil I virkelig lade en pige komme imellem jer?

\scene Der er stille lidt mens Kalkulus og Lille blå kigger ned i gulvet. 

\says{B} Nej... det vil vi ikke. 

\says{L} I to er BEDSTE venner! 

\scene Kalkulus og Lille blå kigger hinanden i øjnene. 

\says{K} Vi ER bedste venner! Vi følges altid ad. 

\scene Kalkulus og Lille Blå bliver mere og mere rørte.

\says{B} Så mange timer sammen i diverse skoletasker.. 

\says{K} Så mange sene aftener med frustrerede rus.. 

\says{B} Du er den bedste ven man kan have! 

\says{K} Nej DU er den bedste ven man kan have! Og vet du hvad? Vi kan godt deles om Anna

\scene{Bøgerne virker næsten glade, da en telefon pludselig ringer}

\says{A} "Hej det Anna"

\scene{Bøgerne stopper pludselig forskrækket}

\says{A} "Ja jeg skal nok komme senere til ZFC kampen, jeg skal bare lige mødes med en rus og sælge min kalkulus"

\scene{Kalkulus bliver enormt ked af det, de andre bøger kigger bedrøvet på Kalkulus}

\says{A} "Nej hvem har dog også brug for sådan en dum, tung, norsk bog... Vi ses!"

\says{K}\act{Kaster sig på knæ, og kaster armene i vejret} Neeeeeiiiiii!

\scene Lys ned

\end{sketch}
\end{document}